\chapter{Transformations}\label{geometry:ch:transformations}

\begin{center}\itshape
Given a manifoldness and a group of transformations of the same; to investigate the configurations belonging to the
manifoldness with regard to such properties as are not altered by the transformations of the group.
\upshape(Felix Klein, the Erlangen programme, 1872, in the translation of 1893)
\end{center}

In October 1872 Felix Klein, aged twenty-three, became a full professor at the small university of Erlangen in Bavaria. The custom
was to give an inaugural lecture and to publish a programme of research, and Klein's pamphlet, \emph{Comparative considerations on recent
geometric research}, changed the meaning of the word geometry. In the nineteenth century geometries had multiplied: Euclidean,
affine, projective, spherical, hyperbolic, and nobody knew how they were related. Klein's answer was that a geometry is a group. Choose
the transformations you allow; the geometry is the study of whatever they leave unchanged. Euclid allows rigid motions and studies
lengths and angles; the projective geometer allows the transformations of perspective and studies only incidence and the cross-ratio.
A geometry with more transformations sees fewer properties, and the geometries are ordered like their groups.

Long before Klein, artists had studied the groups of their patterns without naming them. The Moorish craftsmen of the Alhambra in
Granada covered its walls with periodic tilings of astonishing variety. In 1891 the Russian crystallographer Evgraf Fedorov proved that
there are exactly seventeen kinds of symmetry for a periodic pattern of the plane, and in 1924 George P\'olya rediscovered the count
and drew the seventeen patterns. The Dutch artist Maurits Escher, who had copied the Alhambra tiles in 1922 and 1936, read P\'olya's
paper and used it to design his famous tessellations of birds, fish and lizards.

This chapter is about the transformations of the plane that preserve distances (isometries) and shapes (similarities), the groups
they form, and the patterns those groups classify. We prove that every isometry is a product of at most three reflections, that a
finite symmetry group of a figure is cyclic or dihedral, that there are exactly seven kinds of friezes, and that no periodic pattern
has five-fold symmetry. Groups are as in Volume~2, Chapters~1 to~4.

\begin{think}
A \emph{frieze} is a pattern along a strip that repeats by a translation, like a decorated border on a vase or a row of footprints in
the sand. How many essentially different kinds of symmetry can a frieze have? Think about footprints of someone walking, someone
hopping on one foot, and someone jumping with both feet together.
\end{think}

\section{Isometries}

An \emph{isometry} of $\R^n$ is a map $f$ with $|f(x)-f(y)|=|x-y|$ for all $x,y$. Translations, rotations and reflections are examples.

\begin{theorem}[isometries are affine]\label{geometry:thm:isometryaffine}
Every isometry of $\R^n$ has the form $f(x)=Ax+b$ with $A$ an orthogonal matrix and $b\in\R^n$. In particular it is a bijection, and
the isometries form a group.
\end{theorem}
\begin{proof}
Put $b=f(0)$ and $g(x)=f(x)-b$, an isometry with $g(0)=0$. Then $|g(x)|=|x|$ for all $x$, and from
$|x-y|^2=|x|^2+|y|^2-2\langle x,y\rangle$ applied to $x,y$ and to $g(x),g(y)$ we get $\langle g(x),g(y)\rangle=\langle x,y\rangle$. Now
expand
\[
|g(x+y)-g(x)-g(y)|^2
\]
into nine inner products. Each of them equals the corresponding inner product without $g$, so the whole expression equals
$|(x+y)-x-y|^2=0$, and $g(x+y)=g(x)+g(y)$. In the same way $|g(\lambda x)-\lambda g(x)|^2=|\lambda x-\lambda x|^2=0$. So $g$ is linear and
preserves inner products: its matrix $A$ satisfies $A^{\mathsf T}A=I$. An orthogonal matrix is invertible, so $f$ is a bijection, and
the inverse $x\mapsto A^{\mathsf T}(x-b)$ is again an isometry.
\end{proof}

In the plane, identified with $\C$, an orthogonal matrix has orthonormal columns: the first is $(\cos\theta,\sin\theta)$ for some $\theta$, and
the second, a unit vector perpendicular to it, is $\pm(-\sin\theta,\cos\theta)$. The sign $+$ gives the rotation $z\mapsto e^{i\theta}z$ and the sign $-$
gives the reflection $z\mapsto e^{i\theta}\bar z$. So every isometry of the plane is
\begin{equation}\label{geometry:eq:isomC}
\begin{gathered}
z\mapsto e^{i\theta}z+b\quad\text{(direct)}\qquad\text{or}\\
z\mapsto e^{i\theta}\bar z+b\quad\text{(indirect)},
\end{gathered}
\end{equation}
according to whether it preserves or reverses orientation.

\begin{theorem}[three reflections suffice]\label{geometry:thm:threereflections}
An isometry of the plane is determined by the images of three non-collinear points, and it is a product of at most three
reflections in lines.
\end{theorem}
\begin{proof}
A point $P$ is determined by its distances to three non-collinear points $A,B,C$: the points at given distances from $A$ and $B$ are the
intersections of two circles, at most two points, symmetric about the line $AB$; their distances to $C$ differ, because $C$ is not on
the line $AB$ which is their perpendicular bisector. So if two isometries agree on $A,B,C$, they agree everywhere.

Now let $f$ be an isometry and $A,B,C$ non-collinear. If $f(A)\ne A$, let $s_1$ be the reflection in the perpendicular bisector of
$A$ and $f(A)$, so that $s_1(A)=f(A)$; otherwise $s_1=\mathrm{id}$. Let $B_1=s_1(B)$. If $B_1\ne f(B)$, let $s_2$ be the reflection in the
perpendicular bisector of $B_1$ and $f(B)$; otherwise $s_2=\mathrm{id}$. The point $f(A)$ is on that bisector, because
$|f(A)-B_1|=|s_1(A)-s_1(B)|=|A-B|=|f(A)-f(B)|$, so $s_2$ fixes $f(A)$, and $s_2s_1$ maps $A,B$ to $f(A),f(B)$. In the same way a third
reflection $s_3$ (or the identity) fixes $f(A)$ and $f(B)$ and moves $s_2s_1(C)$ to $f(C)$. The isometries $f$ and $s_3s_2s_1$ agree on
$A,B,C$, so they are equal.
\end{proof}

\begin{figure}[ht]
\centering
\fig{geometry-ch02-tworeflections}
\caption{A reflection in $m_1$ followed by a reflection in $m_2$ is the rotation about their intersection by twice the angle
$\varphi$ from $m_1$ to $m_2$ (Exercise~\ref{geometry:ex:chasles}).}
\label{geometry:fig:tworeflections}
\end{figure}

Composing two reflections in lines meeting at an angle $\varphi$ gives a rotation by $2\varphi$ (Figure~\ref{geometry:fig:tworeflections}), and
two reflections in parallel lines give a translation. So a product of two reflections is direct and one of three is indirect, and
the direct isometries are products of two reflections. The next theorem names all isometries.

\begin{theorem}[classification of plane isometries]\label{geometry:thm:chasles}
Every isometry of the plane is exactly one of the following: the identity; a translation; a rotation about a point by an angle
$\theta\not\equiv0$; a reflection in a line; a \emph{glide reflection}, that is a reflection in a line followed by a non-zero translation
along that line.
\end{theorem}
\begin{proof}
Direct case, $f(z)=e^{i\theta}z+b$. If $e^{i\theta}=1$, $f$ is the translation by $b$ (the identity if $b=0$). Otherwise $f$ has the unique
fixed point $c=b/(1-e^{i\theta})$, and $f(z)-c=e^{i\theta}(z-c)$: it is the rotation about $c$ by $\theta$.

Indirect case, $f(z)=e^{i\theta}\bar z+b$. Conjugate by the rotation $r(z)=e^{i\theta/2}z$:
\[
r^{-1}fr(z)=e^{-i\theta/2}\bigl(e^{i\theta}e^{-i\theta/2}\bar z+b\bigr)=\bar z+b',\qquad b'=e^{-i\theta/2}b.
\]
Write $b'=t+iu$. Then $\bar z+b'=(x+t)+i(u-y)$ for $z=x+iy$: this is the reflection in the horizontal line $y=u/2$ followed by the
translation by $t$ along it, a reflection if $t=0$ and a glide reflection if $t\ne0$. Conjugating back by $r$ gives the same kind of map
for $f$, with the line rotated by $\theta/2$. The five kinds are distinct: they differ in orientation and in their sets of fixed points
(everything, nothing, a point, a line, nothing), and a translation is direct while a glide reflection is not.
\end{proof}

A \emph{similarity} multiplies all distances by the same factor $k>0$; by Theorem~\ref{geometry:thm:isometryaffine} applied to $f/k$, the
similarities of the plane are $z\mapsto az+b$ and $z\mapsto a\bar z+b$ with $a\ne0$. A direct similarity that is not a translation has
exactly one fixed point $c=b/(1-a)$, around which it is a rotation followed by a scaling: a \emph{spiral similarity}.

\section{Klein's Erlangen programme}

Here is Klein's idea in a table. Each group contains the one above it, and preserves less.

\medskip
\begin{center}
\begin{tabular}{@{}p{0.44\linewidth}p{0.5\linewidth}@{}}
\hline
group & what it preserves\\
\hline
isometries $x\mapsto Ax+b$, $A$ orthogonal & lengths, angles, areas\\
similarities $x\mapsto kAx+b$, $k>0$ & angles, ratios of lengths, circles\\
affine maps $x\mapsto Mx+b$, $M$ invertible & lines, parallelism, ratios along a line, ratios of areas\\
projective maps (Appendix~D) & lines, the cross-ratio\\
\hline
\end{tabular}
\end{center}
\medskip

For the affine row: an affine map sends the line $p+tu$ to the line $Mp+b+tMu$, so it preserves lines and parallelism, and it
preserves the parameter $t$, hence ratios along a line; it multiplies all areas by $|\det M|$ (Volume~2, Appendix~B). The table explains why
some theorems are easier than they look. A statement that only involves lines, parallelism and ratios along lines, such as ``the
medians of a triangle meet at a point two thirds of the way down each median'', is affine; and the affine group is so large that it
can move any triangle to any other.

\begin{proposition}\label{geometry:prop:affinetriangles}
For any two triangles $ABC$ and $A'B'C'$ there is exactly one affine map sending $A,B,C$ to $A',B',C'$.
\end{proposition}
\begin{proof}
An affine map with $A\mapsto A'$ is $x\mapsto A'+M(x-A)$. It sends $B$ and $C$ to $B'$ and $C'$ exactly when $M(B-A)=B'-A'$ and
$M(C-A)=C'-A'$. Since $B-A,C-A$ is a basis of $\R^2$, there is exactly one such linear map $M$, and it is invertible because
$B'-A',C'-A'$ is also a basis.
\end{proof}

So an affine statement about triangles needs to be checked only on one triangle, the equilateral one, or the one with vertices $(0,0)$,
$(1,0)$, $(0,1)$: this is how we solved Routh's problem in Exercise~\ref{geometry:ex:routh}. In the same way, every ellipse is the affine image of a
circle, and an affine statement about ellipses (for example, that the midpoints of parallel chords lie on a line through the centre)
need only be checked on a circle.

\section{Finite symmetry groups}

The \emph{symmetry group} of a figure $X$ in the plane is the group of isometries $f$ with $f(X)=X$. A regular $n$-gon has the
\emph{dihedral group} $D_n$ of order $2n$: $n$ rotations about its centre and $n$ reflections. A three-legged
\emph{triskelion} or a pinwheel has only the $n$ rotations, the \emph{cyclic group} $C_n$. Leonardo da Vinci, designing chapels with
radial symmetry, noted that these were the only possibilities, and Hermann Weyl named the theorem after him.

\begin{theorem}[Leonardo's theorem]\label{geometry:thm:leonardo}
A finite group $G$ of isometries of the plane fixes a point, and it is either a cyclic group $C_n$ of rotations about that point, or a
dihedral group $D_n$, made of $n$ rotations and $n$ reflections in lines through that point.
\end{theorem}
\begin{proof}
Take any point $p$ and its orbit $\{g(p):g\in G\}$, with centroid $c=\frac1{|G|}\sum_{g\in G}g(p)$. An isometry $h$ is affine, so it maps
the centroid of a finite set to the centroid of the image, and $h(c)=\frac1{|G|}\sum_gh(g(p))=c$, because $g\mapsto hg$ is a bijection of $G$.
So every element of $G$ fixes $c$; put $c=0$. Then $G$ contains no translations or glide reflections (they fix no point), so its elements
are rotations about $0$ and reflections in lines through $0$.

The rotations form a subgroup $G^+$ (the kernel of the orientation homomorphism $G\to\{\pm1\}$). It is finite; if $G^+\ne\{1\}$, let $\theta$
be the smallest positive angle of a rotation in $G^+$. Every angle $\psi$ of an element of $G^+$ is a multiple of $\theta$: otherwise
$\psi=k\theta+\eta$ with $0<\eta<\theta$, and the rotation by $\eta=\psi-k\theta$ is in $G^+$, a contradiction. Since $2\pi$ is such an
angle, $\theta=2\pi/n$ and $G^+=C_n$. If $G=G^+$ we are done. Otherwise $G^+$ has index $2$ and $G=G^+\cup sG^+$ for any reflection
$s\in G$: these are $n$ rotations and $n$ reflections, and $G$ is the symmetry group of a regular $n$-gon with a vertex on the mirror of $s$,
that is $D_n$.
\end{proof}

\section{Friezes}

\begin{definition}\label{geometry:def:frieze}
A \emph{frieze group} is a group $G$ of isometries of the plane that maps the real axis to itself, and whose translations are exactly the
translations by multiples of a fixed non-zero real number. Scaling, we take that number to be $1$, so that the translations of $G$ form
the group $T=\{z\mapsto z+k:k\in\Z\}$.
\end{definition}

An isometry maps the real axis to itself exactly when, in the form \eqref{geometry:eq:isomC}, $e^{i\theta}=\pm1$ and $b=t$ is real. So the elements
of a frieze group are of four kinds, with $t$ real:
\begin{itemize}
\item $z+t$, a translation;
\item $-z+t$, the half-turn about the point $t/2$;
\item $-\bar z+t$, the reflection in the vertical line $x=t/2$;
\item $\bar z+t$, the reflection in the axis ($t=0$) or a glide reflection along it.
\end{itemize}
Record for each element whether it reverses the direction of the axis and whether it swaps the two sides of the axis: this is a pair of
signs, $(+,+)$, $(-,-)$, $(-,+)$, $(+,-)$ respectively, and composing maps multiplies the signs. We get a homomorphism $\sigma$ from $G$ to the
Klein four-group $V=\{\pm1\}^2$, whose kernel is the set of translations in $G$, that is $T$.

We call two frieze groups of the same \emph{type} if one is conjugate to the other by a similarity mapping the real axis to itself. Such a
conjugation preserves the signs $\sigma$, preserves $T$ when it preserves the unit of length, and maps reflections to reflections.

\begin{theorem}[the seven friezes]\label{geometry:thm:friezes}
There are exactly seven types of frieze groups. Up to type they are generated by $T$ and the following maps:
\begin{center}
\begin{tabular}{@{}lll@{}}
\hline
$\sigma(G)$ & extra generators & name\\
\hline
$\{1\}$ & none & hop\\
$\langle(+,-)\rangle$ & $\bar z+\frac12$ & step\\
$\langle(-,+)\rangle$ & $-\bar z$ & sidle\\
$\langle(-,-)\rangle$ & $-z$ & spinning hop\\
$V$ & $-\bar z$, $\bar z+\frac12$ & spinning sidle\\
$\langle(+,-)\rangle$ & $\bar z$ & jump\\
$V$ & $-\bar z$, $\bar z$ & spinning jump\\
\hline
\end{tabular}
\end{center}
The names are John Conway's. The hop has only translations; the step has glide reflections; the sidle vertical mirrors; the spinning
hop half-turns; the spinning sidle vertical mirrors, glides and half-turns; the jump the horizontal mirror; the spinning jump all of
these.
\end{theorem}
\begin{proof}
The image $\sigma(G)$ is a subgroup of $V$, and $V$ has five subgroups: the trivial one, the three of order $2$, and $V$ itself. Since the
kernel of $\sigma$ is $T$, the group $G$ is the union of the cosets $g_vT$, one for each $v\in\sigma(G)$, where $g_v$ is any element with
$\sigma(g_v)=v$. Conjugating by a translation $z\mapsto z+s$ along the axis changes $-z+t$ into $-z+t-2s$ and $-\bar z+t$ into
$-\bar z+t-2s$, but leaves $\bar z+t$ unchanged.

\emph{$\sigma(G)=\{1\}$.} Then $G=T$: the hop.

\emph{$\sigma(G)=\{(+,+),(-,-)\}$.} Then $G$ contains a half-turn $-z+t$; conjugating with $s=t/2$ we may assume it is $-z$, and
$G=T\cup(-z)T$: the spinning hop.

\emph{$\sigma(G)=\{(+,+),(-,+)\}$.} In the same way we may assume $-\bar z\in G$, and $G=T\cup(-\bar z)T$: the sidle.

\emph{$\sigma(G)=\{(+,+),(+,-)\}$.} Then $G$ contains some $g=\bar z+t$, and $g^2(z)=\overline{\bar z+t}+t=z+2t$ is a translation of $G$, so
$2t\in\Z$. Multiplying $g$ by a translation of $T$ changes $t$ by an integer, so we may take $t=0$ or $t=\frac12$, and the two cases
are distinguished by whether $G$ contains the reflection in the axis (all elements $\bar z+t$ of $G$ have $t\in\Z$ in the first case,
$t\in\frac12+\Z$ in the second). These are the jump and the step.

\emph{$\sigma(G)=V$.} We may assume $-\bar z\in G$ as before, and then $G$ contains some $\bar z+t$ with $t\in\{0,\frac12\}$, and these two
cases are again distinguished by the presence of the horizontal mirror; conjugating by translations does not change $t$. They are the
spinning jump and the spinning sidle. In the spinning sidle, the half-turns are the products $-\bar z\circ(\bar z+\frac12)=-z-\frac12$
and their translates: the centres of the half-turns lie halfway between the mirrors.

Finally each of the seven sets is a group with translation subgroup exactly $T$ (one checks closure on the coset representatives,
as for the spinning sidle above), and two groups in different rows have different $\sigma(G)$ or differ in containing the horizontal
mirror, which a conjugation of the allowed kind preserves. So the seven types are distinct.
\end{proof}

\begin{figure}[ht]
\centering
\fig{geometry-ch02-friezes}
\caption{The seven frieze patterns, each obtained by applying a frieze group to one L-shaped tile. The dotted line is the axis. The
step is the trail of a walker, the jump that of someone jumping with feet together.}
\label{geometry:fig:friezes}
\end{figure}

John Conway's names, chosen with the footprints of the opening question in mind, make the seven types easy to remember: a person who
\emph{hops} on one foot, who \emph{steps} (left, right, left), who \emph{jumps} with both feet together, who \emph{sidles} sideways; and
the \emph{spinning} versions, where the trail is also unchanged by a half-turn. Footprints are only the beginning: the borders of Greek
vases, Navajo blankets and wrought-iron railings all belong to one of the seven types, and anthropologists use the types to compare
the decorative traditions of different cultures.

\section{Lattices and the crystallographic restriction}

A \emph{lattice} in $\R^n$ is a subgroup $\Lambda=\Z v_1\oplus\dots\oplus\Z v_n$ generated by a basis of $\R^n$. The translations of a
periodic pattern of the plane (a wallpaper) form a lattice, the pattern of a crystal has a lattice of translations in $\R^3$, and
Appendix~I studies lattices in their own right.

\begin{think}
Many periodic patterns have symmetries of order $2$, $3$, $4$ or $6$: think of bricks, honeycombs, square tiles. Can you design a
periodic tiling of the floor with a rotation of order $5$ among its symmetries? Of order $8$?
\end{think}

\begin{theorem}[the crystallographic restriction]\label{geometry:thm:crystallographic}
Let $\Lambda$ be a lattice in $\R^2$ or $\R^3$, and let $\rho$ be a rotation of finite order $n$ with $\rho(\Lambda)=\Lambda$. Then
$n\in\{1,2,3,4,6\}$.
\end{theorem}
\begin{proof}
In a basis of $\Lambda$, the linear map $\rho$ sends basis vectors to integer combinations of basis vectors, so its matrix has integer
entries, and its trace is an integer. The trace does not depend on the basis. In an orthonormal basis adapted to the rotation by
$\theta=2\pi k/n$ (with $\gcd(k,n)=1$), the trace is $2\cos\theta$ in the plane and $1+2\cos\theta$ in space. Either way
$2\cos\theta\in\Z$, so $\cos\theta\in\{0,\pm\frac12,\pm1\}$, and $\theta$ is a multiple of $\frac\pi2$ or $\frac\pi3$: the order $n$ is $1$,
$2$, $3$, $4$ or $6$.
\end{proof}

\begin{remark}
Quasicrystals, discovered by Dan Shechtman in 1982 (Nobel Prize in Chemistry, 2011), have diffraction patterns with ten-fold symmetry. They
do not contradict the theorem: they are ordered but not periodic, like Penrose's tilings of the plane by two kinds of rhombi, which have
five-fold symmetry and no lattice of translations.
\end{remark}

The \emph{point group} of a lattice $\Lambda\subset\R^2$ is the group of isometries fixing $0$ and mapping $\Lambda$ to itself. It is
finite: choose $r$ so large that the vectors of $\Lambda$ of length at most $r$ span $\R^2$; there are finitely many of them, the
point group permutes them, and a linear map is determined by what it does to a spanning set. It contains $-1$. By
Theorem~\ref{geometry:thm:leonardo} and Theorem~\ref{geometry:thm:crystallographic} it is $C_n$ or $D_n$ with $n\in\{2,4,6\}$ (it contains the half-turn,
so $n$ is even).

\begin{proposition}[the five lattice types]\label{geometry:prop:fivelattices}
The point group of a plane lattice is one of $C_2$, $D_2$, $D_4$, $D_6$. If it contains a reflection $s$, then $\Lambda$ has a basis
either of the form $a,b$ with $a$ on the mirror of $s$ and $b$ perpendicular to it (a \emph{rectangular} lattice), or of the form
$a,\frac12(a+b)$ with $a,b$ as before (a \emph{centred rectangular} lattice). With the square lattice ($D_4$) and the hexagonal one ($D_6$)
this gives the five types of Figure~\ref{geometry:fig:lattices}.
\end{proposition}
\begin{proof}
For the first claim it remains to exclude $C_4$ and $C_6$: if a rotation of order $4$ or $6$ preserves $\Lambda$, take a vector $v$ of
minimal length; then $\Lambda=\Z v+\Z\rho v$ (Exercise~\ref{geometry:ex:hexagonal} for order $6$, and the same argument with squares instead of triangles for
order $4$), and this lattice is also preserved by the reflection in the line of $v$, which fixes $v$ and maps $\rho v$ to
$\rho^{-1}v$, an element of $\Lambda$. So $G=D_4$ or $D_6$.

Now let $s$ be a reflection preserving $\Lambda$, with mirror line $\ell$. For $v\in\Lambda$, the vector $v+s(v)\in\Lambda$ lies on $\ell$
and $v-s(v)\in\Lambda$ lies on the perpendicular line $\ell^\perp$. These are not all zero, so $\Lambda\cap\ell$ and $\Lambda\cap\ell^\perp$
are non-zero discrete subgroups of lines, generated by vectors $a$ and $b$. The sublattice $\Lambda_0=\Z a+\Z b$ contains $2v=(v+sv)+(v-sv)$
for every $v\in\Lambda$, so $\Lambda\subseteq\frac12\Lambda_0$, and $\Lambda/\Lambda_0$ is a subgroup of $\frac12\Lambda_0/\Lambda_0$, whose
non-zero elements are represented by $\frac a2$, $\frac b2$, $\frac{a+b}2$. But $\frac a2\notin\Lambda$, because it lies on $\ell$ and is
shorter than $a$, and likewise $\frac b2\notin\Lambda$. So either $\Lambda=\Lambda_0$, or $\Lambda=\Lambda_0+\Z\frac{a+b}2$, which has the
basis $a,\frac{a+b}2$.
\end{proof}

\begin{figure}[ht]
\centering
\fig{geometry-ch02-lattices}
\caption{The five types of plane lattices, with point groups $C_2$, $D_2$, $D_2$, $D_4$, $D_6$. In the centred rectangular lattice the
rectangle spanned by $a$ and $b$ has a lattice point at its centre.}
\label{geometry:fig:lattices}
\end{figure}

\begin{remark}
A \emph{wallpaper group} is a group of isometries of the plane whose translations form a lattice. Combining the five lattice types with
the possible point groups and with glide reflections, one finds exactly seventeen wallpaper groups, as Fedorov proved in 1891. The proof
follows the same pattern as Theorem~\ref{geometry:thm:friezes}, with more cases; it is given in full in Armstrong's \emph{Groups and Symmetry}, and Conway's orbifold notation in \emph{The Symmetries of Things} explains the number $17$ by a beautiful counting rule. In
three dimensions there are $230$ space groups, found independently by Fedorov and Sch\"onflies in 1891.
\end{remark}

\section{A gem: the shortest closed billiard path}

Here is a problem of Giovanni Fagnano from 1775 that reflections solve in a few lines. Given an acute triangle, find the triangle of
smallest perimeter with one vertex on each side.

\begin{theorem}[Fagnano]\label{geometry:thm:fagnano}
In an acute triangle $ABC$, the inscribed triangle of minimal perimeter is the \emph{orthic triangle}, whose vertices are the feet of the
three altitudes.
\end{theorem}
\begin{proof}
Fix a point $D$ on $BC$, and let $D'$ and $D''$ be its mirror images in the lines $AB$ and $AC$. For $F$ on $AB$ and $E$ on $AC$,
$DF=D'F$ and $DE=D''E$, so the perimeter of $DEF$ is $D'F+FE+ED''\ge D'D''$, with equality when $F$ and $E$ are the intersections of the
segment $D'D''$ with $AB$ and $AC$ (the triangle being acute, these lie on the sides). Now $AD'=AD=AD''$, and the angle $D'AD''$ is
$2\angle BAC$ (the reflections double the angles $DAB$ and $DAC$), which is less than $180^\circ$. So $D'D''=2\,AD\sin A$, which is smallest
when $AD$ is smallest, that is when $D$ is the foot of the altitude from $A$.

So the minimal perimeter is $2h_a\sin A$, where $h_a$ is the altitude, and it is attained by a unique triangle, the one built from the
foot $D$ of that altitude. Since $2h_a\sin A=2\,AB\sin B\sin A$ is symmetric in the vertices ($h_a=AB\sin B$ and $2R\sin C=AB$ give
$2h_a\sin A=4R\sin A\sin B\sin C$), the same reasoning from $B$ and from $C$ produces the same unique minimiser, whose vertices are therefore
the feet of all three altitudes.
\end{proof}

A billiard ball shot along the orthic triangle would bounce off each side and come back to its start: at each vertex the two sides make
equal angles with the side of $ABC$, because the segment $D'D''$ is straight and reflection preserves angles. Light, which takes the
shortest path, does the same; the reflection trick is Heron's proof of the law of reflection, from the first century.

\begin{beauty}
Seven friezes, seventeen wallpapers, two hundred and thirty crystals: the craftsmen of the Alhambra could not have used an eighteenth
wallpaper symmetry however hard they tried, and no crystal can have the symmetry of a starfish. Klein's insight that a geometry is a
group turns these facts into theorems about a few signs and integer traces, and the same insight will tell us in Appendix~D
what a projective theorem is, and in Appendix~J what the hyperbolic plane is.
\end{beauty}

\section{Exercises}

\begin{exercise}[$\star$]\label{geometry:ex:chasles}
Show that the reflection in a line $m_1$ followed by the reflection in a line $m_2$ is the rotation about $m_1\cap m_2$ by twice the angle
from $m_1$ to $m_2$, if the lines meet, and a translation by twice the vector from $m_1$ to $m_2$, if they are parallel.
\end{exercise}

\begin{exercise}[$\star$]\label{geometry:ex:tworotations}
Show that the rotation by $\alpha$ about $p$ followed by the rotation by $\beta$ about $q$ is a rotation by $\alpha+\beta$ if
$\alpha+\beta\not\equiv0\pmod{2\pi}$, and a translation otherwise. Where is the centre? (Write each rotation as two reflections, one of them in
the line $pq$.)
\end{exercise}

\begin{exercise}[$\star$]\label{geometry:ex:dihedral}
Show that the symmetry group of a regular $n$-gon has order $2n$, that its rotations form a normal subgroup, and that $srs=r^{-1}$ for a
reflection $s$ and a rotation $r$ in it.
\end{exercise}

\begin{exercise}[$\star\star$]\label{geometry:ex:conjugacy}
In the group of isometries of the plane, show that: two rotations are conjugate if and only if their angles are equal or opposite; two
translations are conjugate if and only if their vectors have the same length; all reflections are conjugate; two glide reflections are
conjugate if and only if their translation parts have the same length.
\end{exercise}

\begin{exercise}[$\star$]\label{geometry:ex:affineellipse}
Using Proposition~\ref{geometry:prop:affinetriangles} and the fact that an ellipse is the image of a circle by an affine map, show that the midpoints
of a family of parallel chords of an ellipse lie on a line through its centre.
\end{exercise}

\begin{exercise}[$\star$]\label{geometry:ex:friezeid}
Find the frieze type of each of the infinite words
$\dots\mathsf{LLLL}\dots$, $\dots\mathsf{NNNN}\dots$, $\dots\mathsf{VVVV}\dots$, $\dots\mathsf{DDDD}\dots$, $\dots\mathsf{HHHH}\dots$,
$\dots\mathsf{pbpb}\dots$ and $\dots\textsf{VΛVΛ}\dots$ (letters equally spaced, the axis through their middle).
\end{exercise}

\begin{exercise}[$\star$]\label{geometry:ex:fivefold}
Show without the trace that a lattice in the plane has no rotation of order $5$: take two lattice points at minimal distance and
rotate.
\end{exercise}

\begin{exercise}[$\star\star$]\label{geometry:ex:hexagonal}
Let $\Lambda$ be a plane lattice preserved by the rotation $\rho$ by $60^\circ$, and let $v$ be a non-zero vector of $\Lambda$ of minimal
length. Show that $\Lambda=\Z v+\Z\rho v$. Deduce that every such lattice is similar to the Eisenstein integers $\Z[\omega]$.
\end{exercise}

\begin{exercise}[$\star\star$]\label{geometry:ex:cube}
Show that the rotation group of the cube has order $24$, in two ways: by counting the images of a face with a marked edge, and by listing the
rotations by their axes.
\end{exercise}

\begin{exercise}[$\star\star$]\label{geometry:ex:billiards}
A billiard ball moves in the unit square, bouncing off the sides with equal angles. Show, by reflecting the square instead of the path,
that the trajectory with initial direction of slope $p/q$ (in lowest terms) is periodic, and find its length.
\end{exercise}

\begin{exercise}[$\star$]\label{geometry:ex:glidesquare}
Show that the square of a glide reflection is a translation, and that a glide reflection is determined by its square up to choosing the
axis among the lines parallel to the translation. Is every translation the square of a glide reflection?
\end{exercise}

\begin{exercise}[$\star\star\star$]\label{geometry:ex:finiteSO3}
Let $G$ be a finite group of rotations of $\R^3$ fixing $0$, of order $N>1$. Each non-identity element fixes two \emph{poles} on the unit
sphere. By counting pairs (element, pole) in two ways, show that
\[
2\Bigl(1-\frac1N\Bigr)=\sum_{\text{orbits of poles}}\Bigl(1-\frac1{n_i}\Bigr),
\]
where $n_i$ is the order of the stabiliser of a pole in the $i$-th orbit, and deduce that there are two or three orbits, with
$(n_i)$ equal to $(n,n)$, $(2,2,n)$, $(2,3,3)$, $(2,3,4)$ or $(2,3,5)$.
\end{exercise}

\section{Hints}

\begin{hint}{geometry:ex:chasles}
Choose coordinates with $m_1$ the real axis and use the formula $z\mapsto e^{2i\varphi}\bar z$ for the reflection in a line through $0$
at angle $\varphi$.
\end{hint}
\begin{hint}{geometry:ex:tworotations}
Write the first rotation as $s_2s_1$ with $s_2$ the reflection in $pq$, and the second as $s_3s_2$.
\end{hint}
\begin{hint}{geometry:ex:dihedral}
Use Theorem~\ref{geometry:thm:leonardo} and the orientation homomorphism.
\end{hint}
\begin{hint}{geometry:ex:conjugacy}
Conjugating an isometry by $h$ moves its fixed points, axes and translation vectors by $h$.
\end{hint}
\begin{hint}{geometry:ex:affineellipse}
The statement is affine, so check it for a circle.
\end{hint}
\begin{hint}{geometry:ex:friezeid}
Look for vertical mirrors, a horizontal mirror, half-turns and glides.
\end{hint}
\begin{hint}{geometry:ex:fivefold}
If $\rho$ has order $5$ and $u,v\in\Lambda$ are at minimal distance, look at $\rho^2$ applied around a point, or at $v+\rho^2v$.
\end{hint}
\begin{hint}{geometry:ex:hexagonal}
The points of $\Z v+\Z\rho v$ cut the plane into equilateral triangles of side $|v|$. Every point of the plane is at distance less than
$|v|$ from one of their vertices.
\end{hint}
\begin{hint}{geometry:ex:cube}
The stabiliser of a face with a marked edge is trivial.
\end{hint}
\begin{hint}{geometry:ex:billiards}
Unfold: the reflected squares tile the plane, and the path becomes a straight line.
\end{hint}
\begin{hint}{geometry:ex:glidesquare}
Use the normal form $z\mapsto\bar z+t$.
\end{hint}
\begin{hint}{geometry:ex:finiteSO3}
Each non-identity element has exactly $2$ poles; a pole with stabiliser of order $n_i$ has an orbit of $N/n_i$ poles. Then
$\sum(1-1/n_i)<2$ and each term is at least $\frac12$.
\end{hint}

\section{Solutions}

\begin{solution}{geometry:ex:chasles}
\idea Reflections are $z\mapsto e^{2i\varphi}\bar z$ through the origin; composing two of them cancels the conjugations.
\writeup If the lines meet, put the intersection at $0$ and $m_1$ along the real axis. Then $s_1(z)=\bar z$ and $s_2(z)=e^{2i\varphi}\bar z$, so
$s_2s_1(z)=e^{2i\varphi}\overline{\bar z}=e^{2i\varphi}z$, the rotation by $2\varphi$. If they are parallel, take $m_1$ the real axis and $m_2$
the line $\operatorname{Im}z=d$; then $s_2(z)=\bar z+2id$ and $s_2s_1(z)=z+2id$, the translation by twice the vector $id$ from $m_1$ to $m_2$.
\end{solution}

\begin{solution}{geometry:ex:tworotations}
\idea Write rotations as pairs of reflections so that the middle ones cancel.
\writeup Let $\ell$ be the line $pq$. By Exercise~\ref{geometry:ex:chasles}, the rotation by $\alpha$ about $p$ is $s_\ell s_1$ where $s_1$ is the reflection in the
line through $p$ making the angle $-\alpha/2$ with $\ell$, and the rotation by $\beta$ about $q$ is $s_2s_\ell$ where $s_2$ is the reflection in the
line through $q$ at angle $\beta/2$ with $\ell$. The composite is $s_2s_\ell s_\ell s_1=s_2s_1$. The lines of $s_1$ and $s_2$ make the angle
$\frac{\alpha+\beta}2$; if it is not a multiple of $\pi$ they meet at a point $c$, and $s_2s_1$ is the rotation about $c$ by $\alpha+\beta$. The
centre $c$ is the third vertex of the triangle with base $pq$ and base angles $\alpha/2$ at $p$ and $\beta/2$ at $q$. Otherwise the lines are
parallel and the composite is a translation.
\end{solution}

\begin{solution}{geometry:ex:dihedral}
\idea The rotations are the kernel of the orientation map to $\{\pm1\}$.
\writeup By Theorem~\ref{geometry:thm:leonardo} the group is $C_m$ or $D_m$; it contains the $n$ rotations by multiples of $2\pi/n$ and the reflection in a
line through the centre and a vertex, and no rotation by a smaller angle (it would not map vertices to vertices). So it is $D_n$, of
order $2n$. The rotations are the kernel of the homomorphism to $\{\pm1\}$ given by orientation, hence a normal subgroup. Finally $srs$ is
direct, and in complex coordinates with $s(z)=\bar z$ and $r(z)=e^{i\theta}z$, $srs(z)=\overline{e^{i\theta}\bar{\bar z}}=e^{-i\theta}z=r^{-1}(z)$.
\end{solution}

\begin{solution}{geometry:ex:conjugacy}
\idea Conjugation is a change of coordinates: $hgh^{-1}$ does to $h(x)$ what $g$ does to $x$.
\writeup If $g$ is the rotation by $\theta$ about $c$ and $h$ is direct, $hgh^{-1}$ is the rotation by $\theta$ about $h(c)$; if $h$ is indirect, the
angle becomes $-\theta$ (in coordinates, conjugating $z\mapsto e^{i\theta}z$ by $z\mapsto\bar z$ gives $z\mapsto e^{-i\theta}z$). Since all points are images of one another under translations, the class of the rotation is determined by $\pm\theta$. If
$g$ is the translation by $v$, $hgh^{-1}$ is the translation by $Av$ where $A$ is the linear part of $h$, and $|Av|=|v|$; conversely any two
vectors of equal length are related by a rotation. A reflection in $\ell$ is conjugated to the reflection in $h(\ell)$, and any line is the
image of any other. A glide reflection along $\ell$ with translation $v$ becomes the glide along $h(\ell)$ with translation $Av$; the length
$|v|$ is an invariant (the square is the translation by $2v$), and two glides with $|v|=|v'|$ are conjugate by an isometry taking the first
axis to the second and $v$ to $v'$.
\end{solution}

\begin{solution}{geometry:ex:affineellipse}
\idea Affine maps preserve lines, parallelism, midpoints and centres of symmetry, so we may replace the ellipse by a circle.
\writeup For a circle, the midpoint of a chord is the foot of the perpendicular from the centre, so the midpoints of chords perpendicular to a
direction $u$ lie on the diameter perpendicular to $u$. An affine map $f$ taking the circle to the ellipse maps parallel chords to parallel
chords, midpoints to midpoints and that diameter to a line through the centre of the ellipse (the image of the centre of the circle).
\end{solution}

\begin{solution}{geometry:ex:friezeid}
\idea For each word, list which of the four kinds of symmetry occur and read Theorem~\ref{geometry:thm:friezes} backwards.
\writeup $\mathsf{LLLL}$: translations only, a hop. $\mathsf{NNNN}$: the letter $\mathsf N$ has a half-turn symmetry and no mirror, a spinning hop.
$\mathsf{VVVV}$: vertical mirrors through each letter and between letters, no half-turn, a sidle. $\mathsf{DDDD}$: the horizontal mirror, a
jump. $\mathsf{HHHH}$: both mirrors and half-turns, a spinning jump. $\mathsf{pbpb}$: the reflection in the axis maps $\mathsf p$ to $\mathsf b$, so
the glide reflection by half a period is a symmetry, but no mirror is: a step. $\textsf{VΛVΛ}$: vertical mirrors through each
letter, a glide mapping $\mathsf V$ to $\Lambda$, half-turns between them, no horizontal mirror: a spinning sidle.
\end{solution}

\begin{solution}{geometry:ex:fivefold}
\idea A rotation of order $5$ produces two lattice points closer than the minimal distance.
\writeup Let $\delta>0$ be the minimal distance between points of $\Lambda$ and let $\rho$ (rotation by $72^\circ$ about $0$) preserve $\Lambda$.
Take $v\in\Lambda$ with $|v|=\delta$. Then $v+\rho^2v\in\Lambda$, and the angle between $v$ and $\rho^2v$ is $144^\circ$, so
$|v+\rho^2v|=2\delta\cos72^\circ\approx0.618\,\delta$, which is positive and less than $\delta$, a contradiction.
\end{solution}

\begin{solution}{geometry:ex:hexagonal}
\idea A lattice point missing from $\Z v+\Z\rho v$ would be too close to one that is there.
\writeup Let $\Lambda_0=\Z v+\Z\rho v\subseteq\Lambda$. Its points are the vertices of a tiling of the plane by equilateral triangles of side
$\delta=|v|$ (note $\rho^2v=\rho v-v$). If $w\in\Lambda\setminus\Lambda_0$, it lies in one of these triangles, and its distance to the
nearest vertex $u$ is at most the circumradius $\delta/\sqrt3<\delta$. Then $w-u\in\Lambda$ is non-zero and shorter than $v$, a contradiction.
So $\Lambda=\Lambda_0$, and the similarity $z\mapsto vz$ (in complex notation) maps $\Z[\omega]=\Z+\Z e^{i\pi/3}$ onto it.
\end{solution}

\begin{solution}{geometry:ex:cube}
\idea Orbit-stabiliser, or an inventory of axes.
\writeup A rotation that maps a face to itself and one of its edges to itself fixes the centre of the face and the midpoint of the edge,
as well as the centre of the cube; it fixes two independent vectors, so it is the identity. Hence a rotation is determined by the
image of one face with a marked edge. There are $6$ faces with $4$ edges each, and all $24$ choices occur (rotate the face into place,
then turn about its axis), so the order is $24$. By axes: the identity; around each of the $3$
axes through centres of opposite faces, rotations by $90^\circ$, $180^\circ$, $270^\circ$ ($9$ in all); around each of the $4$ diagonals, by
$120^\circ$ and $240^\circ$ ($8$); around each of the $6$ axes through midpoints of opposite edges, by $180^\circ$ ($6$). Total
$1+9+8+6=24$. (The group is isomorphic to $S_4$, acting on the four diagonals.)
\end{solution}

\begin{solution}{geometry:ex:billiards}
\idea Instead of bouncing the ball, reflect the table: the path becomes a straight line through a tiling of the plane by squares.
\writeup Each bounce off a side is replaced by continuing straight into the mirror image of the square in that side. The reflected squares are
the unit squares of the grid $\Z^2$, and folding them back is the map $\R^2\to[0,1]^2$ that reduces each coordinate modulo $2$ and then
folds $[1,2]$ onto $[0,1]$; it is invariant under translations by $2\Z^2$. The straight line from $x_0$ in direction $(q,p)$ reaches
$x_0+(2q,2p)$, a translate by a vector of $2\Z^2$, after length $2\sqrt{p^2+q^2}$, and the ball is back at its starting point with its
starting direction: the trajectory is periodic. It cannot close up earlier: a return to the same point with the same direction
corresponds to a translation by a vector of $2\Z^2$ (an even number of reflections in each direction), and the multiples $t(q,p)$ in
$2\Z^2$ are the multiples of $(2q,2p)$, because $\gcd(p,q)=1$. So the period is $2\sqrt{p^2+q^2}$ (for paths avoiding the corners).
\end{solution}

\begin{solution}{geometry:ex:glidesquare}
\idea In the normal form of Theorem~\ref{geometry:thm:chasles} everything is explicit.
\writeup After an isometric change of coordinates a glide is $g(z)=\bar z+t$ with $t>0$ real, along the real axis; then
$g^2(z)=z+2t$. Conversely, the glides with square $z\mapsto z+2t$ are the maps $z\mapsto\bar z+t+iu$ for real $u$, which are glides along the
horizontal lines $\operatorname{Im}z=u/2$: one for each line parallel to the translation. So every translation is the square of a glide
(indeed of many).
\end{solution}

\begin{solution}{geometry:ex:finiteSO3}
\idea Count each (non-identity element, pole) pair once by elements and once by poles.
\writeup Each of the $N-1$ non-identity elements has $2$ poles, so there are $2(N-1)$ pairs. A pole $x$ is fixed by the $n_x$ elements of its
stabiliser, $n_x-1$ of them non-trivial, and its orbit has $N/n_x$ poles, all with the same stabiliser size. So the count by poles is
$\sum_i\frac N{n_i}(n_i-1)$, and dividing by $N$ gives the formula. The left side is in $[1,2)$ and each term on the right is in $[\frac12,1)$,
so there are $2$ or $3$ orbits. With $2$ orbits, $\frac1{n_1}+\frac1{n_2}=\frac2N$ with $n_i\le N$ forces $n_1=n_2=N$ (cyclic group). With $3$
orbits, $\frac1{n_1}+\frac1{n_2}+\frac1{n_3}=1+\frac2N>1$ with $n_1\le n_2\le n_3$ forces $n_1=2$, then $\frac1{n_2}+\frac1{n_3}>\frac12$ forces
$n_2=2$ (any $n_3$: dihedral) or $n_2=3$ and $n_3\in\{3,4,5\}$ (the rotation groups of the tetrahedron, cube and icosahedron, of orders
$12$, $24$, $60$).
\end{solution}

\section*{Sources and further reading}

\begin{itemize}
\item F.~Klein, \emph{Vergleichende Betrachtungen \"uber neuere geometrische Forschungen}, Erlangen, 1872; English translation \emph{A
comparative review of recent researches in geometry}, Bulletin of the New York Mathematical Society 2 (1893), 215--249.
\item M.~A.~Armstrong, \emph{Groups and Symmetry}, Springer, 1988: isometries, frieze and wallpaper groups, the crystallographic restriction,
finite rotation groups.
\item J.~H.~Conway, H.~Burgiel and C.~Goodman-Strauss, \emph{The Symmetries of Things}, A~K~Peters, 2008: the orbifold notation and the names of
the friezes.
\item H.~Weyl, \emph{Symmetry}, Princeton University Press, 1952: Leonardo's theorem, ornaments and crystals, by a great mathematician writing
for everyone.
\item G.~P\'olya, \emph{\"Uber die Analogie der Kristallsymmetrie in der Ebene}, Zeitschrift f\"ur Kristallographie 60 (1924), 278--282;
D.~Schattschneider, \emph{The plane symmetry groups: their recognition and notation}, American Mathematical Monthly 85 (1978), 439--450.
\item D.~Shechtman, I.~Blech, D.~Gratias and J.~W.~Cahn, \emph{Metallic phase with long-range orientational order and no translational
symmetry}, Physical Review Letters 53 (1984), 1951--1953.
\item T.~Hawkins, \emph{The Erlanger Programm of Felix Klein: reflections on its place in the history of mathematics}, Historia
Mathematica 11 (1984), 442--470 (Klein's appointment and the reception of the programme).
\item D.~Schattschneider, \emph{Visions of Symmetry: Notebooks, Periodic Drawings, and Related Work of M.~C.~Escher}, W.~H.~Freeman, 1990
(Escher's visits to the Alhambra in 1922 and 1936 and his reading of P\'olya).
\item D.~K.~Washburn and D.~W.~Crowe, \emph{Symmetries of Culture: Theory and Practice of Plane Pattern Analysis}, University of Washington
Press, 1988 (frieze types in anthropology).
\item T.~L.~Heath, \emph{A History of Greek Mathematics}, vol.~2, Oxford, 1921 (Heron's proof of the law of reflection in his
\emph{Catoptrica}).
\item H.~S.~M.~Coxeter, \emph{Introduction to Geometry}, 2nd ed., Wiley, 1969, Chapters~3 and~4 (isometries, Fagnano's problem).
\end{itemize}
